% GENERATED FILE. Edit canonical lessons, then run npm run generate:books.
% canonical-source-hash: fnv1a64:97973b04e8acd30d
% canonical-lessons: SA-C503-sahakari, SA-C503-daksah, SA-C503-parisramah, SA-C503-niyokta, SA-C503-niyunkte

\chapter{People at Work}
\label{ch:sa-people-at-work}
\hlchaptermodality{503}

\begin{chapteropening}
\textbf{By the end of this chapter:} \emph{I can talk about the people I work with and what they are good at.}

Complete the last lesson of chapter 503: I can talk about the people I work with and what they are good at.
\end{chapteropening}

\section[\texorpdfstring{\sk{सहकारी} (sahakārī)}{सहकारी (sahakārī)}]{\sk{सहकारी} (sahakārī) — a colleague}

\label{lesson:SA-C503-sahakari}



\begin{quote}
Before the new one: say the Sanskrit for a saw, then the Sanskrit for a chisel.
\end{quote}

\subsection*{You'll want to know: \sk{सहकारी}}
\textbf{\sk{सहकारी}} — \emph{sahakārī} — \textquotedblleft{}a colleague\textquotedblright{}.

\begin{quote}
\sk{मम} \sk{सहकारी} \sk{अपि} \sk{अत्र} \sk{कार्यं} \sk{करोति।}
\end{quote}

\begin{quote}
\emph{mama sahakārī api atra kāryaṃ karoti.}
\end{quote}

\begin{quote}
My colleague works here too.
\end{quote}

A masculine noun.

\begin{cousinweb}
\textbf{\sk{सहकारी}} (\emph{sahakārī}) is built from \emph{saha}, \textquotedblleft{}together\textquotedblright{}, and \emph{kārin}, \textquotedblleft{}doing\textquotedblright{}, from the root √kṛ, \textquotedblleft{}to do\textquotedblright{}; MW gives \emph{sahakārin} as one who acts together with another, an assistant; \textquotedblleft{}a colleague\textquotedblright{} is its use today (MW).
\end{cousinweb}

\begin{grammarlens}[title={What you've built}]
One more word for describing what you lived and felt.
\end{grammarlens}

\subsection*{Your turn}
\begin{itemize}
\raggedright
  \item \emph{Say it:} \emph{sahakārī}
  \item \emph{Say it:} \emph{sahakārī}, once more
  \item \emph{Say it:} \emph{mama sahakārī api atra kāryaṃ karoti.} — My colleague works here too.
  \item \emph{Recall:} say the Sanskrit for a saw, then the Sanskrit for a chisel, then say \emph{sahakārī} again
\end{itemize}

\begin{tcolorbox}[breakable,colback=teal!4,colframe=teal!35!black,title={Before you move on}]
What does \sk{सहकारी} mean? (\textquotedblleft{}A colleague\textquotedblright{}.) Say it once more.
\end{tcolorbox}
\section[\texorpdfstring{\sk{दक्षः} (dakṣaḥ)}{दक्षः (dakṣaḥ)}]{\sk{दक्षः} (dakṣaḥ) — adroit; capable}

\label{lesson:SA-C503-daksah}



\begin{quote}
Before the new one: say the Sanskrit for a chisel, then the Sanskrit for a colleague.
\end{quote}

\subsection*{You'll want to know: \sk{दक्षः}}
\textbf{\sk{दक्षः}} — \emph{dakṣaḥ} — \textquotedblleft{}adroit; capable\textquotedblright{}.

\begin{quote}
\sk{सः} \sk{दक्षः} \sk{तक्षकः} \sk{अस्ति।}
\end{quote}

\begin{quote}
\emph{saḥ dakṣaḥ takṣakaḥ asti.}
\end{quote}

\begin{quote}
He is a capable carpenter.
\end{quote}

\begin{cousinweb}
\textbf{\sk{दक्षः}} (\emph{dakṣaḥ}) comes from the root √dakṣ, \textquotedblleft{}to be able, to be strong\textquotedblright{}, under which MW gives the adjective \emph{dakṣa}, \textquotedblleft{}able, fit, adroit, clever\textquotedblright{} (MW).
\end{cousinweb}

\begin{grammarlens}[title={What you've built}]
One more word for describing what you lived and felt.
\end{grammarlens}

\subsection*{Your turn}
\begin{itemize}
\raggedright
  \item \emph{Say it:} \emph{dakṣaḥ}
  \item \emph{Say it:} \emph{dakṣaḥ}, once more
  \item \emph{Say it:} \emph{saḥ dakṣaḥ takṣakaḥ asti.} — He is a capable carpenter.
  \item \emph{Recall:} say the Sanskrit for a chisel, then the Sanskrit for a colleague, then say \emph{dakṣaḥ} again
\end{itemize}

\begin{tcolorbox}[breakable,colback=teal!4,colframe=teal!35!black,title={Before you move on}]
What does \sk{दक्षः} mean? (\textquotedblleft{}Adroit; capable\textquotedblright{}.) Say it once more.
\end{tcolorbox}
\section[\texorpdfstring{\sk{परिश्रमः} (pariśramaḥ)}{परिश्रमः (pariśramaḥ)}]{\sk{परिश्रमः} (pariśramaḥ) — hard work}

\label{lesson:SA-C503-parisramah}



\begin{quote}
Before the new one: say the Sanskrit for a colleague, then the Sanskrit for adroit.
\end{quote}

\subsection*{You'll want to know: \sk{परिश्रमः}}
\textbf{\sk{परिश्रमः}} — \emph{pariśramaḥ} — \textquotedblleft{}hard work\textquotedblright{}.

\begin{quote}
\sk{परिश्रमेण} \sk{सः} \sk{धनं} \sk{लभते।}
\end{quote}

\begin{quote}
\emph{pariśrameṇa saḥ dhanaṃ labhate.}
\end{quote}

\begin{quote}
He gets money through hard work.
\end{quote}

A masculine noun.

\begin{cousinweb}
\textbf{\sk{परिश्रमः}} (\emph{pariśramaḥ}) is built on \emph{pari}, \textquotedblleft{}all round\textquotedblright{}, and the root √śram, \textquotedblleft{}to toil, to grow weary\textquotedblright{}; MW gives \emph{pariśrama} the senses \textquotedblleft{}exertion, labour\textquotedblright{} (MW).
\end{cousinweb}

\begin{grammarlens}[title={What you've built}]
One more word for describing what you lived and felt.
\end{grammarlens}

\subsection*{Your turn}
\begin{itemize}
\raggedright
  \item \emph{Say it:} \emph{pariśramaḥ}
  \item \emph{Say it:} \emph{pariśramaḥ}, once more
  \item \emph{Say it:} \emph{pariśrameṇa saḥ dhanaṃ labhate.} — He gets money through hard work.
  \item \emph{Recall:} say the Sanskrit for a colleague, then the Sanskrit for adroit, then say \emph{pariśramaḥ} again
\end{itemize}

\begin{tcolorbox}[breakable,colback=teal!4,colframe=teal!35!black,title={Before you move on}]
What does \sk{परिश्रमः} mean? (\textquotedblleft{}Hard work\textquotedblright{}.) Say it once more.
\end{tcolorbox}
\section[\texorpdfstring{\sk{नियोक्ता} (niyoktā)}{नियोक्ता (niyoktā)}]{\sk{नियोक्ता} (niyoktā) — an employer}

\label{lesson:SA-C503-niyokta}



\begin{quote}
Before the new one: say the Sanskrit for adroit, then the Sanskrit for hard work.
\end{quote}

\subsection*{You'll want to know: \sk{नियोक्ता}}
\textbf{\sk{नियोक्ता}} — \emph{niyoktā} — \textquotedblleft{}an employer\textquotedblright{}.

\begin{quote}
\sk{मम} \sk{नियोक्ता} \sk{उदारः} \sk{अस्ति।}
\end{quote}

\begin{quote}
\emph{mama niyoktā udāraḥ asti.}
\end{quote}

\begin{quote}
My employer is generous.
\end{quote}

A masculine noun.

\begin{cousinweb}
\textbf{\sk{नियोक्ता}} (\emph{niyoktā}) is built on \emph{ni} and the root √yuj, \textquotedblleft{}to yoke, to appoint\textquotedblright{}, with the agent ending \emph{-tṛ}; MW gives \textquotedblleft{}one who fastens\textquotedblright{} and \textquotedblleft{}a master\textquotedblright{}; \textquotedblleft{}an employer\textquotedblright{} is the modern sense (MW).
\end{cousinweb}

\begin{grammarlens}[title={What you've built}]
One more word for describing what you lived and felt.
\end{grammarlens}

\subsection*{Your turn}
\begin{itemize}
\raggedright
  \item \emph{Say it:} \emph{niyoktā}
  \item \emph{Say it:} \emph{niyoktā}, once more
  \item \emph{Say it:} \emph{mama niyoktā udāraḥ asti.} — My employer is generous.
  \item \emph{Recall:} say the Sanskrit for adroit, then the Sanskrit for hard work, then say \emph{niyoktā} again
\end{itemize}

\begin{tcolorbox}[breakable,colback=teal!4,colframe=teal!35!black,title={Before you move on}]
What does \sk{नियोक्ता} mean? (\textquotedblleft{}An employer\textquotedblright{}.) Say it once more.
\end{tcolorbox}
\section[\texorpdfstring{\sk{नियुंक्ते} (niyuṅkte)}{नियुंक्ते (niyuṅkte)}]{\sk{नियुंक्ते} (niyuṅkte) — employs; appoints}

\label{lesson:SA-C503-niyunkte}



\begin{quote}
Before the new one: say the Sanskrit for hard work, then the Sanskrit for an employer.
\end{quote}

\subsection*{You'll want to know: \sk{नियुंक्ते}}
\textbf{\sk{नियुंक्ते}} — \emph{niyuṅkte} — \textquotedblleft{}employs; appoints\textquotedblright{}.

\begin{quote}
\sk{नियोक्ता} \sk{मम} \sk{भ्रातरं} \sk{नियुंक्ते।}
\end{quote}

\begin{quote}
\emph{niyoktā mama bhrātaraṃ niyuṅkte.}
\end{quote}

\begin{quote}
The employer takes on my brother.
\end{quote}

Past (he …): \sk{सः} \sk{मम} \sk{भ्रातरं} \sk{नियुक्तवान्}, \emph{saḥ mama bhrātaraṃ niyuktavān}.

\begin{cousinweb}
\textbf{\sk{नियुंक्ते}} (\emph{niyuṅkte}) is built on \emph{ni}, \textquotedblleft{}down, into\textquotedblright{}, and the root √yuj, \textquotedblleft{}to yoke\textquotedblright{}; MW gives \emph{niyuj} the senses \textquotedblleft{}to fasten to\textquotedblright{} and \textquotedblleft{}to appoint, to employ\textquotedblright{}; its past participle is \emph{niyukta} (MW).
\end{cousinweb}

\begin{grammarlens}[title={What you've built}]
One more word for describing what you lived and felt.
\end{grammarlens}

\subsection*{Your turn}
\begin{itemize}
\raggedright
  \item \emph{Say it:} \emph{niyuṅkte}
  \item \emph{Say it:} \emph{niyuṅkte}, once more
  \item \emph{Say it:} \emph{niyoktā mama bhrātaraṃ niyuṅkte.} — The employer takes on my brother.
  \item \emph{Recall:} say the Sanskrit for hard work, then the Sanskrit for an employer, then say \emph{niyuṅkte} again
\end{itemize}

\begin{tcolorbox}[breakable,colback=teal!4,colframe=teal!35!black,title={Before you move on}]
What does \sk{नियुंक्ते} mean? (\textquotedblleft{}Employs; appoints\textquotedblright{}.) Say it once more.
\end{tcolorbox}
